\chapter{Tổng quan và kiến thức nền}
\label{chap:background}

Chương này trình bày bối cảnh kỹ thuật của khóa luận: nhà thông minh dùng cảm biến môi trường
(mục~\ref{sec:bg-smarthome}), quy trình nhận dạng hoạt động (mục~\ref{sec:bg-har}), bài toán nhà
nhiều người với trọng tâm là luận án của Chen~\cite{chen2021thesis} (mục~\ref{sec:bg-multi}),
các hướng khai phá mẫu sinh hoạt và theo dõi thay đổi hành vi (mục~\ref{sec:bg-routine}), xu hướng
dùng mô hình ngôn ngữ (mục~\ref{sec:bg-llm}), các kiến thức nền được dùng trong khóa luận
(mục~\ref{sec:bg-tools}) và cuối cùng là khoảng trống mà khóa luận hướng tới
(mục~\ref{sec:bg-gap}).

\section{Nhà thông minh và cảm biến môi trường}
\label{sec:bg-smarthome}

Theo cách phân loại phổ biến, cảm biến dùng cho nhận dạng hoạt động gồm năm nhóm: camera, thiết bị
đeo, điện thoại, âm thanh và cảm biến môi trường~\cite{chen2021thesis,bouchabou2021survey}. Camera
cho thông tin giàu nhất nhưng đắt và xâm phạm riêng tư; thiết bị đeo rẻ hơn nhưng người cao tuổi
thường ngại đeo hoặc quên đeo. \emph{Cảm biến môi trường} --- gắn cố định trong nhà --- được xem là
hướng hợp lý cho chăm sóc tại nhà vì không xâm nhập, rẻ và dễ lắp đặt~\cite{chen2021thesis}.

Dự án CASAS của Đại học Bang Washington~\cite{cook2013casas} là nguồn dữ liệu chuẩn của lĩnh vực.
Mỗi nhà gắn vài chục cảm biến: cảm biến chuyển động (ký hiệu \texttt{M}, giá trị \texttt{ON/OFF}),
cảm biến cửa (\texttt{D}, \texttt{OPEN/CLOSE}) và cảm biến nhiệt độ (\texttt{T}). Mỗi dòng dữ liệu
là một \emph{sự kiện}: \texttt{2009-06-10 07:12:05 M021 ON}. Một số dòng kèm nhãn do người chú
thích ghi, dạng \texttt{R2\_Take\_Medicine begin/end}, đánh dấu đầu và cuối một lần thực hiện hoạt
động. Đặc điểm quan trọng: dữ liệu \emph{thưa} (cảm biến chỉ kích hoạt khi có chuyển động),
\emph{không đều} theo thời gian và \emph{vô danh} --- không ghi ai đã kích hoạt cảm biến.

\section{Nhận dạng hoạt động từ cảm biến môi trường}
\label{sec:bg-har}

Chen~\cite{chen2021thesis} tổng kết quy trình nhận dạng hoạt động dựa trên dữ liệu thành bốn bước:
\emph{thu thập dữ liệu} $\rightarrow$ \emph{tiền xử lý và phân đoạn} $\rightarrow$ \emph{phân lớp}
$\rightarrow$ \emph{suy luận và đánh giá}.

\textbf{Phân đoạn.} Chuỗi sự kiện liên tục phải được cắt thành các đoạn để phân lớp. Ba cách cơ
bản là cửa sổ trượt cố định theo thời gian, cửa sổ cố định theo số sự kiện và cửa sổ động dựa trên
tương quan giữa cảm biến~\cite{chen2021thesis,krishnan2014streaming}. Krishnan và
Cook~\cite{krishnan2014streaming} đề xuất gán nhãn cho \emph{từng sự kiện} dựa trên cửa sổ các sự
kiện ngay trước nó --- cách làm phù hợp với dữ liệu dạng luồng và được khóa luận kế thừa.
Phát hiện điểm thay đổi (CPD)~\cite{aminikhanghahi2017cpd} là một cách khác để tìm ranh giới hoạt
động.

\textbf{Phân lớp.} Các mô hình truyền thống gồm Naive Bayes, mô hình Markov ẩn, trường ngẫu nhiên
có điều kiện, cây quyết định, SVM và rừng ngẫu nhiên~\cite{breiman2001rf}; các mô hình học sâu gồm
CNN và nhất là các mạng hồi quy LSTM, GRU và biến thể hai chiều~\cite{cho2014gru,schuster1997brnn}.
Thực nghiệm của Chen cho thấy mạng hồi quy thường vượt các mô hình khác trên dữ liệu chuỗi thời
gian~\cite{chen2021thesis}.

\textbf{Đánh giá.} Các chỉ số thường dùng là accuracy, precision, recall và F1. Một điểm cần lưu ý
về phương pháp luận: nhiều công trình, kể cả~\cite{chen2021thesis}, dùng kiểm định chéo 10 lần mà
không nói rõ có chia theo thời gian hay không. Với dữ liệu chuỗi thời gian, nếu chia ngẫu nhiên thì
các cửa sổ liền kề --- gần như giống nhau --- rơi vào cả tập huấn luyện lẫn tập kiểm tra, làm kết
quả cao hơn thực tế. Khóa luận chia dữ liệu \emph{theo ngày, theo thứ tự thời gian} để tránh rủi ro
này.

\section{Nhà nhiều người}
\label{sec:bg-multi}

\subsection{Thách thức}

Khi có từ hai người trở lên, các sự kiện của họ trộn vào một luồng theo thứ tự thời
gian~\cite{benmansour2015survey}. Bên cạnh hoạt động tuần tự, xen kẽ và đồng thời của một người,
xuất hiện thêm hoạt động \emph{song song} (mỗi người làm một việc ở một nơi) và \emph{cộng tác}
(cùng làm một việc)~\cite{chen2021thesis}. Hệ thống không chỉ phải biết ``đang có hoạt động gì'' mà
còn phải biết ``ai đang làm''~\cite{singla2010multiresident}.

\subsection{Luận án của Chen (2021)}

Luận án tiến sĩ của Dong Chen~\cite{chen2021thesis} là tài liệu định hướng chính của khóa luận. Luận
án giải ba bài toán con:
\begin{enumerate}
  \item \textbf{Phân đoạn} (Chương 3 của luận án, công bố trên IEEE Internet of Things
  Journal~\cite{chen2021fcmcpd}): gom cảm biến theo \emph{vị trí} bằng Fuzzy C-means --- hai sự kiện
  liên tiếp ở hai vùng khác nhau khó do cùng một người gây ra --- rồi dùng CPD để tìm điểm chuyển
  hoạt động trong từng nhóm.
  \item \textbf{Phân lớp} (Chương 4): mô hình TRANS-BiGRU kết hợp Transformer với Bi-GRU, so sánh với
  khoảng 15 mô hình học máy và học sâu trên CASAS và ARAS~\cite{alemdar2013aras}.
  \item \textbf{Học trực tuyến} (Chương 5): rừng ngẫu nhiên thích nghi kèm bộ phát hiện trôi khái
  niệm (LED-ARF) cho dữ liệu đến liên tục.
\end{enumerate}
Chính luận án nêu ba hạn chế: dữ liệu chỉ có hai người và hoạt động đơn giản; hiệu năng còn khoảng
cách để cải thiện; và các phương pháp \emph{chỉ nhận ra hoạt động hiện tại}, chưa dự đoán hay phát
hiện bất thường. Khóa luận kế thừa \emph{khung} của luận án (tách người $\rightarrow$ nhận dạng) và
\emph{ý tưởng} dùng vị trí để tách người, nhưng không tái hiện các mô hình của luận án; thay vào đó,
khóa luận đi tiếp theo hướng luận án còn bỏ ngỏ: từ nhận dạng sang \emph{hành vi dài hạn}.

\subsection{Gán sự kiện cho từng người}

Bài toán ``sự kiện này do ai gây ra'' được gọi là \emph{liên kết dữ liệu} (data association) hay
\emph{theo dõi nhiều cư dân}. Crandall và Cook~\cite{crandall2009coping} dùng luật dựa trên đồ thị
kề của cảm biến; Wang và Cook~\cite{wang2021smrt} (sMRT) vector hóa cảm biến để theo dõi nhiều người
mà không cần sơ đồ nhà. Gần đây, Chen, Cumin và cộng sự~\cite{chen2026mdp} mô hình hóa bài toán như
một quá trình quyết định Markov, huấn luyện bằng học bắt chước với Transformer nhỏ hoặc mô hình ngôn
ngữ lớn; công trình này cho thấy gán sai người làm sai các chỉ số sức khỏe phía sau, như thời lượng
hoạt động. Các công trình trên đánh giá ở mức \emph{sự kiện} hoặc \emph{vệt}; chưa công trình nào
đánh giá tác động lên \emph{thói quen}.

\section{Khai phá mẫu sinh hoạt và thay đổi hành vi}
\label{sec:bg-routine}

Một hướng khác không hỏi ``hoạt động hiện tại là gì'' mà hỏi ``sinh hoạt thường ngày như thế nào''.
Duckworth và cộng sự~\cite{duckworth2019unsupervised} cho robot di động tự học các lớp hoạt động
không cần nhãn; bộ dữ liệu dài hạn của dự án STRANDS~\cite{coppola2016temporal} là dữ liệu thật của
đề tài thực tập. DISCOVER~\cite{karpekov2025discover} học biểu diễn tự giám sát rồi phân cụm để tìm
``mẫu sinh hoạt hằng ngày'', chỉ cần gán nhãn rất ít cụm --- nhưng chỉ kiểm chứng trên nhà một người.
X-BCD~\cite{civitarese2026xbcd} phát hiện thay đổi hành vi theo thời gian và giải thích cho bác sĩ.

Đề tài thực tập của tác giả thuộc hướng này: thói quen được định nghĩa là nhóm phiên cùng phòng, bắt
đầu lệch nhau không quá khoảng 30 phút, lặp lại ở ít nhất 50\% số ngày, và được tìm bằng DBSCAN
trên không gian giờ vòng tròn. Hạn chế được ghi nhận khi đó: một hoạt động kéo dài cả buổi (ngồi
phòng khách) bị DBSCAN nối thành một khung giờ rất rộng, và dữ liệu thật chỉ có một người.

\section{Xu hướng mô hình ngôn ngữ và mô hình nền}
\label{sec:bg-llm}

Từ 2024, nhiều công trình chuyển chuỗi cảm biến thành văn bản để dùng mô hình ngôn ngữ: ADL-LLM nhận
dạng hoạt động không cần nhãn huấn luyện~\cite{civitarese2024adlllm}; TDOST mô tả mỗi lần cảm biến
kích hoạt bằng câu chữ để chuyển mô hình sang nhà khác~\cite{thukral2024tdost}; một số công trình
dùng biểu diễn ngôn ngữ thay vì sinh văn bản~\cite{dhekane2025prompt}; DomusFM là mô hình nền đầu
tiên huấn luyện riêng cho dữ liệu nhà thông minh~\cite{fiori2026domusfm}. Bộ dữ liệu MuRAL~\cite{chen2025mural}
(2--4 người, chú thích bằng ngôn ngữ tự nhiên) cho thấy các mô hình ngôn ngữ hiện tại vẫn gặp khó
khi phải phân biệt ai kích hoạt cảm biến. Dự đoán hoạt động tiếp theo bằng mô hình ngôn ngữ cũng đã
được thử nghiệm~\cite{doctorarastoo2026prediction}. Các hướng này nằm \emph{ngoài phạm vi} khóa luận
và được nhắc lại trong phần hướng phát triển.

\section{Kiến thức nền}
\label{sec:bg-tools}

\subsection{Thống kê vòng tròn}
Giờ trong ngày là đại lượng tuần hoàn: 23:50 và 00:10 chỉ cách nhau 20 phút nhưng nếu biểu diễn bằng
số thực (23,83 và 0,17) thì cách nhau gần một ngày. Cách xử lý chuẩn~\cite{fisher1993circular} là
đặt giờ $h$ lên vòng tròn với góc $\theta = 2\pi h/24$. Trung bình vòng tròn của các giờ
$h_1,\dots,h_n$ là
\begin{equation}
  \bar h = \frac{24}{2\pi}\,\operatorname{atan2}\Big(\tfrac1n\textstyle\sum_i \sin\theta_i,\;
  \tfrac1n\sum_i \cos\theta_i\Big) \bmod 24,
\end{equation}
và chênh lệch có dấu giữa hai giờ là $\delta(a,b) = \big((a-b+12) \bmod 24\big) - 12 \in [-12, 12)$.

\subsection{Phân cụm theo mật độ}
DBSCAN~\cite{ester1996dbscan} gom các điểm có ít nhất \texttt{min\_samples} láng giềng trong bán
kính $\varepsilon$ thành cụm, các điểm còn lại là nhiễu; không cần biết trước số cụm. Nhược điểm
liên quan trực tiếp tới khóa luận: các điểm có thể ``nối chuỗi'' qua nhau, nên một hoạt động diễn
ra rải rác cả ngày tạo thành một cụm kéo dài cả ngày. HDBSCAN~\cite{campello2013hdbscan} xây cây
phân cấp mật độ và chọn cụm ổn định nhất, không cần $\varepsilon$ nhưng vì thế cũng không có khái
niệm ``lệch không quá 30 phút''.

\subsection{Rừng ngẫu nhiên và mạng hồi quy}
Rừng ngẫu nhiên~\cite{breiman2001rf} kết hợp nhiều cây quyết định học trên các mẫu bootstrap và tập
con đặc trưng ngẫu nhiên; ổn định, ít cần tinh chỉnh và hỗ trợ trực tiếp bài toán nhiều đầu ra. GRU
\cite{cho2014gru} là mạng hồi quy có cổng cập nhật và cổng đặt lại; phiên bản hai chiều
\cite{schuster1997brnn} đọc chuỗi theo cả hai chiều rồi ghép trạng thái ẩn.

\subsection{Phân cụm phổ}
Phân cụm phổ~\cite{vonluxburg2007spectral} dùng các vector riêng của ma trận Laplacian chuẩn hóa
$L = I - D^{-1/2} A D^{-1/2}$ của đồ thị tương đồng $A$. Số cụm có thể chọn theo \emph{khoảng cách
trị riêng}: $k$ là vị trí có bước nhảy lớn nhất giữa hai trị riêng liên tiếp nhỏ nhất.

\subsection{Bài toán ghép cặp}
Để so khớp một--một hai tập phần tử theo một ma trận chi phí, thuật toán Hungary~\cite{kuhn1955hungarian}
tìm cách ghép có tổng chi phí nhỏ nhất trong thời gian đa thức.

\section{Khoảng trống nghiên cứu và vị trí của khóa luận}
\label{sec:bg-gap}

Bảng~\ref{tab:gap} đặt khóa luận cạnh các công trình gần nhất.

\begin{table}[htbp]
\centering
\caption{Vị trí của khóa luận so với các công trình liên quan}
\label{tab:gap}
\small
\begin{tabularx}{\textwidth}{>{\raggedright\arraybackslash}X ccc}
\toprule
Công trình & Nhà nhiều người & Đầu ra là thói quen & Đánh giá mức thói quen \\
\midrule
Chen 2021~\cite{chen2021thesis,chen2021fcmcpd} & \checkmark & -- & -- \\
sMRT 2021~\cite{wang2021smrt}; MDP 2026~\cite{chen2026mdp} & \checkmark & -- & -- \\
DISCOVER 2025~\cite{karpekov2025discover} & -- & \checkmark & một phần \\
X-BCD 2026~\cite{civitarese2026xbcd} & -- & thay đổi hành vi & -- \\
Đề tài thực tập & -- & \checkmark & \checkmark \\
\textbf{Khóa luận này} & \checkmark & \checkmark & \checkmark \\
\bottomrule
\end{tabularx}
\end{table}

Các công trình về nhà nhiều người dừng ở mức sự kiện hoặc hoạt động; các công trình về mẫu sinh hoạt
chỉ làm với một người. Khóa luận đứng ở giao điểm: \emph{học thói quen của từng người khi luồng cảm
biến bị trộn}, đánh giá trực tiếp ở mức thói quen, và đo xem lỗi ở khâu nhận dạng --- trong đó có
lỗi nhầm người --- ảnh hưởng tới thói quen đến đâu.
